\section{The original adelic chart, rational counts and uniform Sobolev bound}
\label{sec:adelic-sobolev}

We reconstruct the actual fixed-level Sobolev producer at source commit
\[
 \text{\ttfamily 7cbf7a0c32fca1badf05931879a37c3f8a2b03b4}.
\]
The formal source is assumed correct.  This chapter gives independent
human proofs and asserts no Lean replay.
Its final assertion is conditional on the original target
\(\mathrm{FinPair}\); neither that finite-volume datum nor boundedness
of an initial theta integral is produced here.

Fix the original \(D,\theta,V\), with \(\mathrm{IsInfFam}(\theta)\).
Write \(F=D.\mathrm{host}.F\), \(E=D.\mathrm{host}.E\),
\(n=V.\mathrm{rank}\), \(I=D.\mathrm{what}_2.\mathrm{Tau}\),
\(G=\mathrm{GA}(D,V)\), and \(\Gamma=\mathrm{GQ}(V)\).
Put \(H_\tau=\mathrm{gramAt}(D,V,\tau)\).
The actual real component of \(g\) is \(g_\tau=\mathrm{archM}(V,\theta_\tau,g)\).
We use exactly the source height
\[
 \mathcal N(g)=\mathrm{normInf}(\theta,V,g)
   =\sum_{\tau\in I}\sum_{a,b<n}
       \bigl(|(g_\tau)_{ab}|+|(g_\tau^{-1})_{ab}|\bigr).
 \tag{ASB.1}
\]
The absolute values are complex norms; no \(1+\mathcal N\), equivalent
height, or rescaled measure is substituted.

For a pair \(p=(\tau,X)\), \(X^*H_\tau+H_\tau X=0\), write
\[
 c_p(t)=\mathrm{oneP}(\theta,V,\tau,X,t),\qquad
 R_p\psi(g)=\left.\frac{d}{dt}\right|_{0}\psi(gc_p(t)).
 \tag{ASB.2}
\]
For a list \(l=[p_1,\ldots,p_k]\), \(R_l\) is the original right fold
\(R_{p_1}\cdots R_{p_k}\), with the rightmost operator acting first.
The Hilbert derivative word \(d_l\) has the same convention.

% BEGIN ASB1
\paragraph{ASB1: the actual real coordinates and one-place Lie basis.}
\label{claim:ASB1}
The real group of the original target is
\[
 G_{\mathbb R}=\prod_{\tau\in I}U(H_\tau),\qquad
 \mathfrak g=\prod_{\tau\in I}
       \{X:X^*H_\tau+H_\tau X=0\}.
 \tag{ASB.3}
\]
There is a real basis \(Y_1,\ldots,Y_r\) of \(\mathfrak g\), each
supported at one place: \(Y_j=\mathrm{single}_{\tau_j}X_j\).
The ordered product
\[
 \Theta_{\mathbb R}(t)=\prod_{j=1}^r e^{t_jY_j},
 \qquad
 \Theta(t)=\prod_{j=1}^r c_{(\tau_j,X_j)}(t_j)
 \tag{ASB.4}
\]
uses the displayed order.  It consists of the actual original
one-parameter subgroups and is continuous.  The actual lift
\(L:G_{\mathbb R}\to G_\infty\) satisfies
\(L(\Theta_{\mathbb R}(t))=\Theta(t)\).

\emph{Proof.}
HCU14 identifies the original tensor product at infinity with its
matrix components, proves that every \(H_\tau\) is Hermitian
nonsingular, and constructs the actual continuous real lift \(L\).
It preserves products on \(G_{\mathbb R}\), is inverse to the real
component map on \(G_\infty\), and agrees with the literal
\(\mathrm{atTau}\) on one-place matrices.
These are identities on the original group, not an abstract
isomorphism chosen for this chapter.

Multiplication by \(H_\tau\) identifies its Lie space with the
skew-Hermitian matrices, since
\((H_\tau X)^*=X^*H_\tau=-H_\tau X\).
It therefore has a real basis, of size \(n^2\).  Take a basis in each
place and concatenate them; these vectors are independent and span
the direct product.  Thus \(r=|I|n^2\), although only finiteness is
needed.  Differentiating
\(e^{tX^*}H_\tau e^{tX}\) gives zero, so \(e^{tX}\in U(H_\tau)\).
The exponential identities for a fixed matrix give
\(c_p(0)=1\) and \(c_p(s+t)=c_p(s)c_p(t)\); continuity follows
from the matrix exponential and the actual lift.
The product identity in (ASB.4) follows from the homomorphism law of
\(L\).  HCU14 also gives the continuous finite part
\[
 a(g)=gL((g_\tau)_\tau)^{-1}\in G_f,
 \qquad g=a(g)L((g_\tau)_\tau),
 \tag{ASB.5}
\]
and commutation of every finite-supported and real-supported element.
RQC1 supplies the original locally compact, second-countable group
topology and Borel structure used below.
% END ASB1

% BEGIN ASB2
\paragraph{ASB2: uniformly controlled shifted second-kind coordinates.}
\label{claim:ASB2}
For the same ordered basis there are \(s>0\), \(R>0\), \(K\ge1\),
and a neighborhood \(\mathcal V\) of \(1\) in \(G_{\mathbb R}\)
such that, for each \(v\in\mathcal V\), one map
\(f_v:\mathbb R^r\to\mathbb R^r\) satisfies, on \(Q=[0,s]^r\),
\[
 \|f_v(t)\|_\infty\le R,\qquad
 \Theta_{\mathbb R}(f_v(t))v=\Theta_{\mathbb R}(t),\qquad
 \|t-t'\|_\infty\le K\|f_v(t)-f_v(t')\|_\infty.
 \tag{ASB.6}
\]
The constants precede \(v,t,t'\).

\emph{Proof.}
We construct the necessary group chart explicitly.  Near zero in
\(\mathfrak g\), the Cayley map
\[
 \mathcal C(X)=(1+X)(1-X)^{-1}
 \tag{ASB.7}
\]
has values in \(G_{\mathbb R}\).  Indeed
\((1+X)^*H(1+X)=(1-X)^*H(1-X)\), place by place.
Its inverse near \(1\) is \((v-1)(v+1)^{-1}\); substituting
\(v^*Hv=H\) proves that this inverse is in \(\mathfrak g\).
These maps are smooth because matrix inversion is a rational map
where its determinant is nonzero.

In this chart the derivative at zero of
\(t\mapsto\Theta_{\mathbb R}(t)\) is
\(\frac12\sum_jt_jY_j\), an isomorphism.
The finite-dimensional inverse function theorem therefore makes
\(\Theta_{\mathbb R}\) a smooth bijection from a neighborhood of
zero onto a group neighborhood, with smooth inverse \(b\).
Choose concentric closed sup-norm balls, the largest still inside
this coordinate neighborhood.  By continuity and compactness,
after shrinking a neighborhood of \(1\), both
\[
 f_v(t)=b(\Theta_{\mathbb R}(t)v^{-1}),\qquad
 g_v(z)=b(\Theta_{\mathbb R}(z)v)
 \tag{ASB.8}
\]
are defined on the required smaller balls, uniformly in \(v\).
Choose \(Q\) sufficiently small that all \(f_v(Q)\) lie in one
closed ball \(B_R\), while \(g_v\) is defined on a neighborhood of
the larger closed ball \(B_{2R}\).
These choices are possible because \(f_1(0)=0\).

The derivative of \(g_v\) in \(z\) is continuous jointly in \(v,z\).
Restrict \(v\) to a neighborhood with compact closure inside the
previous one.  Compactness of that closure times \(B_{2R}\)
gives one operator-norm bound \(K\ge1\).
The line segment between \(f_v(t)\) and \(f_v(t')\) lies in
\(B_R\).  The mean-value integral along the segment and
\(g_v(f_v(t))=t\) prove the final inequality in (ASB.6).
The first two assertions follow from (ASB.8).
Extend \(f_v\) arbitrarily outside its coordinate domain; the
source contract only restricts it on \(Q\).
For \(r=0\), the real group and parameter space are singletons.
Take \(s=R=K=1\), the unique map, and the whole group as
\(\mathcal V\).  Every assertion then holds literally.
% END ASB2

% BEGIN ASB3
\paragraph{ASB3: the full original adelic chart on one open level.}
\label{claim:ASB3}
For every original open subgroup \(U\le G\), the preceding
\(r,\tau_j,X_j,s,R,K\) and a neighborhood \(W\) of \(1\) in \(G\)
can be chosen so that for every \(x\in W\) there are
one \(a_x\in U\cap G_f\) and one \(f_x\) satisfying
\[
 \begin{gathered}
 \|\!f_x(t)\!\|_\infty\le R,\qquad
             \Theta(f_x(t))x=\Theta(t)a_x\quad(t\in Q),\\
 \|t-t'\|_\infty\le K\|f_x(t)-f_x(t')\|_\infty
                       \quad(t,t'\in Q).
 \end{gathered}
 \tag{ASB.9}
\]
In particular \(a_x\) is independent of \(t\).  This is the entire
original \texttt{SobLoc.exists\_adelic\_chart} conclusion.

\emph{Proof.}
Let \(\operatorname{ar}(x)=(x_\tau)_\tau\) be the actual real
component map and let \(a(x)\) be (ASB.5).
Both are continuous, \(\operatorname{ar}(1)=1\), and \(a(1)=1\).
Set
\[
 W=\operatorname{ar}^{-1}(\mathcal V)\cap a^{-1}(U).
 \tag{ASB.10}
\]
This is a neighborhood of \(1\).  For \(x\in W\), take
\(a_x=a(x)\), \(v=\operatorname{ar}(x)\), and \(f_x=f_v\).
The finite part belongs to \(U\cap G_f\).
Lift the middle identity of (ASB.6) through \(L\), and commute
the finite part past \(\Theta(f_x(t))\).  Using
\(x=a_xL(v)\), the result is precisely
\(\Theta(f_x(t))x=\Theta(t)a_x\).
The two bounds are unchanged.  The matrices \(X_j\), ordered
products and one-place lifts are exactly those of ASB1.
No subgroup replacing the supplied \(U\) has been introduced.
% END ASB3

% BEGIN ASB4
\paragraph{ASB4: the scalar cube-to-Haar comparison.}
\label{claim:ASB4}
For the fixed original \(D,\theta,V,U\) and any original left Haar
measure \(\mu\) on \(G\), there are a compact \(\Omega\subset G\)
and \(C_{\mathrm{ch}}\ge0\), chosen before the function, such that
every continuous nonnegative \(h:G\to\mathbb R\) satisfying
\(h(ga)=h(g)\) for \(a\in U\cap G_f\) obeys
\[
          \int_Q h(\Theta(t))\,dt
              \le C_{\mathrm{ch}}\int_\Omega h(g)\,d\mu(g).
 \tag{ASB.11}
\]
This is the actual scalar comparison required by the original
smooth-representative construction.

\emph{Proof.}
First recall, with a proof, the finite-dimensional distortion
estimate being used.  A \(K\)-Lipschitz map between subsets of
\(\mathbb R^r\), for the sup norm, increases Lebesgue outer
measure by at most \(K^r\).  Cover an open set containing its
domain subset by a disjoint dyadic-cube decomposition.  The image
of the portion in a cube of side \(d\) has coordinate ranges of
length at most \(Kd\), so it is covered, with arbitrarily small
enlargement, by a box of volume at most \(K^rd^r\).
Sum these bounds and take the infimum over open covers.
This proves the estimate even for a map defined only on a subset.
For \(r=0\) both spaces are singletons and the same assertion
uses \(K^0=1\).

Consequently, if \(f:A\to B_R\) satisfies the inverse Lipschitz
inequality in (ASB.9), and \(h_1(t)\le h_2(f(t))\) with both
functions nonnegative, then
\[
                   \int_A h_1(t)\,dt
                         \le K^r\int_{B_R}h_2(z)\,dz.
 \tag{ASB.12}
\]
Apply the outer-measure estimate to the inverse of \(f\) on the
image of each superlevel set of \(h_1\).  That image is contained
in the corresponding superlevel set of \(h_2\) in \(B_R\).
Integrating the level-set measures proves (ASB.12) by Tonelli.
No measurability of a separately chosen extension of \(f\) is
needed.

Choose a compact neighborhood \(B\) of \(1\) contained in the
neighborhood \(W\) of ASB3.  Local compactness from RQC1 gives
such a \(B\).  Haar positivity and compact finiteness give
\(0<\mu(B)<\infty\).
For each \(x\in B\), (ASB.9) and the level invariance of \(h\)
give \(h(\Theta(t))=h(\Theta(f_x(t))x)\).
Thus (ASB.12) gives
\[
       A:=\int_Qh(\Theta(t))\,dt
                  \le K^r\int_{B_R}h(\Theta(z)x)\,dz .
 \tag{ASB.13}
\]
Integrate this inequality over \(x\in B\).
Every integrand is continuous on the relevant compact products,
so Fubini applies.  Left Haar invariance gives
\[
 \begin{split}
 \mu(B)A
 &\le K^r\int_{B_R}\int_{\Theta(z)B}h(g)\,d\mu(g)\,dz\\
 &\le K^r\operatorname{vol}(B_R)
                       \int_{\Theta(B_R)B}h(g)\,d\mu(g).
 \end{split}
 \tag{ASB.14}
\]
Take \(\Omega=\Theta(B_R)B\), a compact set, and
\(C_{\mathrm{ch}}=K^r\operatorname{vol}(B_R)/\mu(B)\).
Only left Haar invariance was used; neither a quotient finite
volume nor right Haar invariance is required for this claim.
% END ASB4

% BEGIN ASB5
\paragraph{ASB5: the ordered-word endpoint Sobolev inequality.}
\label{claim:ASB5}
Let \(c_1,\ldots,c_r\) be continuous one-parameter subgroups of
any topological group.  Suppose continuous functions \(F_l\),
indexed by finite words in \(1,\ldots,r\), satisfy
\[
 \left.\frac{d}{dt}\right|_0F_l(gc_j(t))=F_{j::l}(g).
 \tag{ASB.15}
\]
For \(s>0\), if
\[
 \int_{[0,s]^r}|F_l(yc_1(t_1)\cdots c_r(t_r))|^2\,dt\le M
                   \quad(|l|\le r),
 \tag{ASB.16}
\]
then
\[
               |F_{[]}(y)|^2\le(2/s+2s)^rM.
 \tag{ASB.17}
\]
No commutation of the curves or their derivatives is assumed.

\emph{Proof.}
For a continuously differentiable complex function \(f\) on
\([0,s]\), the fundamental theorem and Cauchy--Schwarz give,
for every \(t\in[0,s]\),
\[
 |f(0)|^2\le2|f(t)|^2+2s\int_0^s|f'(u)|^2\,du.
 \tag{ASB.18}
\]
Integrating in \(t\) and dividing by \(s\) proves the
one-dimensional bound with coefficients \(2/s\) and \(2s\).
For \(f(t)=F_l(gc_j(t))\), (ASB.15) and the subgroup law
give its derivative at every \(t\), namely
\(F_{j::l}(gc_j(t))\).  Continuity supplies the fundamental
theorem hypotheses.

Induct on the number of curves.  In dimension zero, the cube
has measure one and (ASB.16) is already the conclusion.
For \(r>0\), apply (ASB.18) in the last curve \(c_r\) to
every word \(l\) in the first \(r-1\) letters with \(|l|\le r-1\),
at \(g=yc_1(t_1)\cdots c_{r-1}(t_{r-1})\).
After integration in the first \(r-1\) variables, the two
full-cube integrals are those of \(F_l\) and \(F_{r::l}\).
They are bounded by \(M\), so the shorter-cube integral is
at most \((2/s+2s)M\).
Apply the induction hypothesis to the first \(r-1\) curves.
This gives (ASB.17), and explains why the derivative is
prepended to the word rather than moved past its other letters.
Our constant is a valid direct choice for the source's
existential constant in
\texttt{sobolev\_iterated\_one\_parameter\_subgroups\_words}.
% END ASB5

% BEGIN ASB6
\paragraph{ASB6: the complete original local Sobolev inequality.}
\label{claim:ASB6}
For the same fixed \(D,\theta,V,U,\mu\), there exist one finite
list \(\mathcal S\) of admissible original Lie words, one compact
\(\Omega\), and \(C_0\ge0\), such that
\[
 |\psi(y)|^2\le C_0\sum_{l\in\mathcal S}
                    \int_\Omega|R_l\psi(yo)|^2\,d\mu(o)
                       \qquad(y\in G).
 \tag{ASB.19}
\]
This holds for every continuous \(\psi\) which is right invariant
under \(U\cap G_f\), is smooth along every original ordered
product of one-place curves, and has continuous \(R_l\psi\) for
every admissible word.  These are precisely the original
\texttt{SobLoc.local\_sobolev} hypotheses.

\emph{Proof.}
Take the ordered basis and cube of ASB3 and the compact set and
constant of ASB4.  List all words of length at most \(r\) in the
finite alphabet \(1,\ldots,r\), and replace each letter \(j\)
by \(p_j=(\tau_j,X_j)\).  This gives one finite list
\(\mathcal S\), with its original right-fold order.
Set \(F_l=R_{p_l}\psi\).

The assumed curve-product smoothness makes the function on
each real slice \(u\mapsto\psi(gL(u))\) smooth in the
second-kind chart of ASB2, and hence in its translates.
Differentiation along \(u\mapsto ue^{tY_j}\) is a smooth vector
field in those coordinates.  Its iterates remain smooth.
Thus (ASB.15) holds for these \(F_l\), with the actual
derivatives of (ASB.2).  Their continuity on all of \(G\)
is exactly the separately assumed word-continuity condition.

Every \(a\in U\cap G_f\) commutes with each \(c_p(t)\), by ASB1.
Induction on the word, differentiating the equal functions,
therefore gives \(F_l(ga)=F_l(g)\).
For fixed \(y\), apply ASB4 to the continuous nonnegative
function \(g\mapsto|F_l(yg)|^2\).
Each cube integral in (ASB.16) is at most
\[
 M=C_{\mathrm{ch}}\sum_{l\in\mathcal S}
                             \int_\Omega|R_l\psi(yo)|^2\,d\mu(o).
 \tag{ASB.20}
\]
All summands are nonnegative.  ASB5 now proves (ASB.19) with
\(C_0=(2/s+2s)^rC_{\mathrm{ch}}\).
All these choices precede \(\psi\) and \(y\).
If \(r=0\), \(\mathcal S\) consists of the empty word and the
same proof applies with the zero-dimensional cube.
% END ASB6

% BEGIN ASB7
\paragraph{ASB7: one denominator for a finite-adelic compact set.}
\label{claim:ASB7}
For a number field \(F_0\) and compact
\(K\subset\mathbb A_{F_0,f}\), there is a positive integer \(m\)
such that every \(a\in F_0\) whose diagonal finite adele belongs
to \(K\) has \(ma\in\mathcal O_{F_0}\).

\emph{Proof.}
The product \(\widehat{\mathcal O}_{F_0}\) of local integer
rings is an open additive subgroup of the finite adeles.
For any finite adele \(z\), only finitely many coordinates are
nonintegral.  A sufficiently large power of each rational prime
below one of these places clears its negative valuation.
The product of those rational prime powers is a positive integer
\(m_z\) with \(m_zz\in\widehat{\mathcal O}_{F_0}\).

Cover \(K\) by finitely many cosets
\(z_j+\widehat{\mathcal O}_{F_0}\), and take \(m=\prod_jm_{z_j}\).
Multiplication by an integer preserves every local integer ring,
so \(mK\subset\widehat{\mathcal O}_{F_0}\).
A rational element integral at all finite places is an algebraic
integer.  To see this directly, take its monic minimal polynomial
over \(\mathbb Q\).  For each rational prime \(p\), each root in
\(\overline{\mathbb Q}_p\) is the image of the element under an
embedding of its generated field; that embedding extends to
\(F_0\), since the extension is algebraic and the target is
algebraically closed.  The resulting \(p\)-adic absolute value
on \(F_0\) is one of its finite places.  All these roots therefore
have \(p\)-adic absolute value at most one.  The elementary
symmetric polynomials in the roots show that every rational
coefficient of the minimal polynomial lies in \(\mathbb Z_p\)
for every \(p\).  A rational number with that property has no
prime in its reduced denominator, so is an integer.
This proves the claimed intersection with the global field.
The argument gives the source's
\texttt{exists\_den\_of\_isCompact}, indeed a denominator for
the whole compact adelic set before restricting to rational points.
% END ASB7

% BEGIN ASB8
\paragraph{ASB8: the exact original tensor-product denominator.}
\label{claim:ASB8}
Let \(E_0/F_0\) be number fields and give
\(\mathbb A_{F_0,f}\otimes_{F_0}E_0\) its original module topology.
For every compact \(K\) in this tensor product there is one
positive integer \(m\) such that
\[
                 1\otimes a\in K\quad\Longrightarrow\quad
                       ma\in\mathcal O_{E_0}.
 \tag{ASB.21}
\]

\emph{Proof.}
Choose an \(F_0\)-basis \(\beta_1,\ldots,\beta_d\) of \(E_0\).
Its tensor coordinates are continuous in the original module
topology, so the union of the finitely many coordinate images of
\(K\) is compact in \(\mathbb A_{F_0,f}\).
ASB7 gives one denominator \(m_1\) for all rational coordinates.
Choose a positive integer \(m_2\) making every \(m_2\beta_j\)
integral over \(\mathbb Z\).
Such an integer exists: if a monic polynomial for \(\beta_j\)
has rational coefficients, clearing their denominators by \(q\)
makes \(q\beta_j\) a root of a monic integer polynomial.

For \(a=\sum_j a_j\beta_j\), the tensor coordinates of
\(1\otimes a\) are exactly the finite diagonal images of \(a_j\).
Thus every \(m_1a_j\) is integral over \(\mathbb Z\), and
\[
                 m_1m_2a=\sum_j(m_1a_j)(m_2\beta_j)
 \tag{ASB.22}
\]
is integral, since algebraic integers are closed under sums
and products and under field inclusion.
Take \(m=m_1m_2\).
This is the full original
\texttt{exists\_den\_of\_isCompact\_tensor}; no comparison with
a differently topologized adelic tensor product is used.
% END ASB8

% BEGIN ASB9
\paragraph{ASB9: integral coordinates bounded by the actual house.}
\label{claim:ASB9}
For a number field \(E_0\) of degree \(d\), write
\(\mathrm{house}(a)=\max_{\sigma:E_0\hookrightarrow\mathbb C}
|\sigma(a)|\).
There is an integral basis \(\omega_1,\ldots,\omega_d\) and
\(c\ge0\) such that every
\(a=\sum_jz_j\omega_j\in\mathcal O_{E_0}\) satisfies
\[
                         |z_j|\le c\,\mathrm{house}(a).
 \tag{ASB.23}
\]
The same conclusion holds for any fixed integral basis, including
the basis stored by the source.

\emph{Proof.}
Here is the algebraic construction behind the coordinate estimate.
Choose a rational basis \(\beta_j\) consisting of algebraic integers
by the denominator argument of ASB8.
The trace Gram matrix \(T_{ij}=\operatorname{Tr}(\beta_i\beta_j)\)
has integer entries and nonzero determinant \(\Delta\).
Nondegeneracy follows directly from separability: in a power basis
of a primitive element the embedding matrix is Vandermonde with
distinct roots, and changing a rational basis preserves
nonsingularity.  A primitive element can be obtained by successively
replacing two generators \(a,b\) with \(a+qb\): only finitely many
rational \(q\) fail to separate their distinct embeddings.

For an algebraic integer \(a=\sum_iq_i\beta_i\), all traces
\(\operatorname{Tr}(a\beta_j)\) are integers.  They are sums of
algebraic integers and are rational; a rational algebraic integer
is an integer by the rational-root test.  Hence
\[
 \sum_j\mathbb Z\beta_j\ \subset\ \mathcal O_{E_0}\
           \subset\ \Delta^{-1}\sum_j\mathbb Z\beta_j.
 \tag{ASB.24}
\]
A subgroup of \(\mathbb Z^d\) is free of rank at most \(d\):
project to the first coordinate, choose a generator of its image,
and apply induction to the kernel.  This also proves finite
generation.  The two inclusions in (ASB.24) therefore give a
free integral basis of rank exactly \(d\).

For any such basis, let \(M_{\sigma j}=\sigma(\omega_j)\).
It is invertible by the same embedding-matrix argument.
The coordinate equation
\[
             (z_j)_j=M^{-1}(\sigma(a))_\sigma
 \tag{ASB.25}
\]
gives (ASB.23) with
\(c=\max_j\sum_\sigma|(M^{-1})_{j\sigma}|\).
The proof uses the actual maximum over all complex embeddings,
so it proves the original \texttt{exists\_repr\_le\_house}
without inserting a lattice-covolume normalization.
% END ASB9

% BEGIN ASB10
\paragraph{ASB10: complete polynomial counting of integral-denominator matrices.}
\label{claim:ASB10}
Put \(d=[E_0:\mathbb Q]\) and \(b=dn^2\).
For every positive integer \(m\) there is \(C_m\ge0\) such that,
for every \(R\ge1\), the set
\[
 \mathcal A(m,R)=
 \{A\in M_n(E_0):mA_{ij}\in\mathcal O_{E_0},\
                         \mathrm{house}(A_{ij})\le R\}
 \tag{ASB.26}
\]
is finite and has at most \(C_mR^b\) elements.
The exponent is chosen before \(m\) and \(R\).

\emph{Proof.}
Use the integral basis and constant \(c\) of ASB9.
Every integral coordinate of every \(mA_{ij}\) has absolute value
at most \(cmR\), because
\(\mathrm{house}(ma)=m\,\mathrm{house}(a)\).
These \(b\) integers determine the matrix injectively, since
\(m\ne0\).  With \(M=\lceil cmR\rceil\), there are at most
\((2M+1)^b\) possible arrays.  Since
\[
             2M+1\le2cmR+3\le(2cm+3)R\quad(R\ge1),
 \tag{ASB.27}
\]
we may take \(C_m=(2cm+3)^b\).
For \(n=0\), the set consists of the unique empty matrix and
the same formula gives exponent zero and constant one.
This proves all quantifiers of the original
\texttt{RatCount.count\_matrices}.
% END ASB10

% BEGIN ASB11
\paragraph{ASB11: the precise original \(\mathrm{normInf}\) estimates.}
\label{claim:ASB11}
For the original group, \(\mathcal N\) is continuous and
nonnegative, and
\[
 \mathcal N(g^{-1})=\mathcal N(g),\qquad
 \mathcal N(gh)\le\mathcal N(g)\mathcal N(h).
 \tag{ASB.28}
\]
If \(n>0\), then \(\mathcal N(g)\ge1\).
If the actual adelic matrix of \(g\) is the diagonal image of
\(A\in M_n(E)\), then
\[
                         \mathrm{house}(A_{ij})\le\mathcal N(g).
 \tag{ASB.29}
\]

\emph{Proof.}
Continuity follows from the actual continuous real component map,
matrix entries and inversion.  Nonnegativity is termwise.
Inversion interchanges the two summands in (ASB.1).
For a complex matrix write \(|A|_1=\sum_{i,j}|A_{ij}|\).
Expanding entries and applying the triangle inequality gives
\(|AB|_1\le|A|_1|B|_1\).
At each place apply this both to \(gh\) and to
\((gh)^{-1}=h^{-1}g^{-1}\), then sum.
The product of the two sums in (ASB.28) contains these
nonnegative terms, as well as additional nonnegative terms.

HCU14's original CM-type dictionary shows that the place set is
nonempty and contains one choice of an extension of every real
embedding of \(F\).
For \(n>0\), choose a place and a diagonal entry of
\(g_\tau g_\tau^{-1}=1\).  It gives
\[
       1\le |g_\tau|_1|g_\tau^{-1}|_1
                 \le\frac14(|g_\tau|_1+|g_\tau^{-1}|_1)^2.
 \tag{ASB.30}
\]
Thus \(\mathcal N(g)\ge2\), which implies the stated bound.
For the house assertion, each complex embedding of the CM field
\(E\) is either the chosen extension at one real place or its
complex conjugate: \(E/F\) is quadratic and \(F\) is totally real.
The absolute values of these two images agree.
The original \(\mathrm{IsInfFam}\) condition evaluates a rational
entry \(A_{ij}\) at the chosen component exactly as that field
embedding.  Each such absolute value is one of the summands in
(ASB.1), proving (ASB.29).
For \(n=0\), \(\mathcal N=0\); no positive-rank assertion is
used in that case.
% END ASB11

% BEGIN ASB12
\paragraph{ASB12: the full quantitative rational-point multiplicity bound.}
\label{claim:ASB12}
For every compact \(\Omega\subset G\), one may take
\[
                         N=2[E:\mathbb Q]\,n^2
 \tag{ASB.31}
\]
so that for every \(x\in G\) there is \(C_x\ge0\) for which,
for every \(u\in G_\infty\) and \(y\in G\), the set
\[
 \mathcal R(x,u,y)=
 \{\gamma\in\Gamma:\exists o\in\Omega,\ \gamma y=xuo\}
 \tag{ASB.32}
\]
is finite and has cardinal at most \(C_x\mathcal N(u)^N\).
This proves
\texttt{RatCount.}\allowbreak\texttt{ncard\_rational\_le}.

\emph{Proof.}
For \(n=0\) the matrix group is a singleton.  Its relevant set
has at most one element; take \(N=0,C_x=1\), with
\(0^0=1\) as in the source.
Suppose \(n>0\), and put \(b=[E:\mathbb Q]n^2\).
The set \(K=\Omega\Omega^{-1}\) is compact.
Choose \(B\ge0\) with \(\mathcal N(w)\le B\) on \(K\).
For fixed \(x\), the finite components of all entries of
\(xwx^{-1}\), \(w\in K\), form a compact subset of the
original finite tensor product.  ASB8 supplies one positive
denominator \(m_x\) for this entire set.  In particular this
choice does not depend on \(u,y\) or a rational point.
Set
\[
                       c_x=\mathcal N(x)^2B+1\ge1.
 \tag{ASB.33}
\]

If (ASB.32) is empty there is nothing to prove.  Otherwise choose
one \(\gamma_0\) and \(o_0\in\Omega\) with
\(\gamma_0y=xuo_0\).  For \(\gamma y=xuo\), cancellation gives
\[
 \delta_\gamma=\gamma\gamma_0^{-1}
                =xu(oo_0^{-1})u^{-1}x^{-1}\in\Gamma.
 \tag{ASB.34}
\]
The map \(\gamma\mapsto\delta_\gamma\) is injective.
Since the finite component of \(u\) is the identity, the finite
component of (ASB.34) is exactly that of
\(x(oo_0^{-1})x^{-1}\).
Therefore the rational matrix \(A_\gamma\) representing
\(\delta_\gamma\) has every \(m_x(A_\gamma)_{ij}\) integral.
By ASB11,
\[
 \begin{split}
 \mathrm{house}((A_\gamma)_{ij})
 &\le\mathcal N(\delta_\gamma)\\
 &\le\mathcal N(x)^2B\,\mathcal N(u)^2
 \le c_x\mathcal N(u)^2 .
 \end{split}
 \tag{ASB.35}
\]
The last radius is at least one by ASB11.
All matrices \(A_\gamma\) consequently lie in
\(\mathcal A(m_x,c_x\mathcal N(u)^2)\).
ASB10 bounds their number by
\[
              C_{m_x}c_x^b\mathcal N(u)^{2b}.
 \tag{ASB.36}
\]
The original matrix realization is injective, so this bounds
(ASB.32).  Take \(C_x=C_{m_x}c_x^b\).
The auxiliary \(\gamma_0,o_0\) only prove an injection; the
constant was chosen without them.  The bound is therefore
uniform in \(y\) and \(u\), with one exponent before \(x\).
% END ASB12

% BEGIN ASB13
\paragraph{ASB13: exact multiplicity-to-\(L^2\) unfolding.}
\label{claim:ASB13}
Fix an original \(\mathrm{FinPair}(V,\mu,\mathcal F)\) and put
\(\nu=(\mu|_{\mathcal F})_*\mathrm{autMk}\).
Let compact \(\Omega\), \(z\in G\) and \(M\) satisfy
\[
          \#\{\gamma\in\Gamma:\gamma y\in z\Omega\}\le M
                       \quad(y\in G),
 \tag{ASB.37}
\]
with each set finite.  If a continuous \(\psi\) represents
\(F_0\in L^2(\Gamma\backslash G,\nu)\) in the original global
almost-everywhere sense, then
\[
                 \int_\Omega|\psi(zo)|^2\,d\mu(o)
                                   \le M^2\|F_0\|_2^2.
 \tag{ASB.38}
\]

\emph{Proof.}
Let \(h(g)=|F_0(\mathrm{autMk}(g))|^2\), using the measurable
representative of the \(L^2\) class.  It is exactly left rational
invariant.  RQC2's countable unfolding, with the indicator of
the compact measurable set \(z\Omega\), gives
\[
 \int_{z\Omega}h(g)\,d\mu(g)
  =\int_{\mathcal F}h(y)
           \#\{\gamma\in\Gamma:\gamma y\in z\Omega\}\,d\mu(y).
 \tag{ASB.39}
\]
For a nonnegative integer \(k\le M\), one has
\(k\le k^2\le M^2\); this also covers \(k=0\).
Thus the right side is at most
\(M^2\int_{\mathcal F}h=M^2\|F_0\|_2^2\), with exactly the
given quotient-measure scale.
The global representative identity gives \(h=|\psi|^2\)
almost everywhere on \(G\), hence on \(z\Omega\).
Left Haar translation turns that integral into the left side of
(ASB.38).  Continuity of \(\psi\) and compactness supply its
ordinary integrability; equivalently the nonnegative argument
already proves finiteness.
This is the original \texttt{integral\_le\_adelic} comparison.
No new fundamental domain or normalized Haar measure is chosen.
% END ASB13

% BEGIN ASB14
\paragraph{ASB14: one exponent and one word list before every vector.}
\label{claim:ASB14}
Fix the original \(D,\theta,V\), \(\mathrm{IsInfFam}(\theta)\),
one \(\mathrm{FinPair}(V,\mu,\mathcal F)\), any original
\(\mathrm{IsReg}\) on its quotient, and one open subgroup \(U\).
There exist \(N\in\mathbb N\) and a finite list \(\mathcal S\)
of admissible original Lie words such that
\[
 \begin{gathered}
 \forall x\in G\ \exists C_x\in\mathbb R\
 \forall F_0,\psi\
 \forall u\in G_\infty,\\
 |\psi(xu)|\le
        C_x\mathcal N(u)^N\sum_{l\in\mathcal S}\|d_lF_0\|_2 .
 \end{gathered}
 \tag{ASB.40}
\]
Here the universally quantified \(F_0,\psi\) satisfy exactly:
\(F_0\) is an original \(\mathrm{IsSmoothVector}'\);
\(\rho(k)F_0=F_0\) for all \(k\in U\cap G_f\);
\(\psi\) is continuous; and \(\mathrm{Repr}(F_0,\psi)\) holds
globally almost everywhere.
This proves the complete original
\texttt{SobolevGrowth.pointwise\_growth} endpoint.

\emph{Proof.}
Apply ASB6 first, obtaining \(\mathcal S,\Omega,C_0\), with
\(C_0\ge0\), independently of \(F_0,\psi,x\).
Apply ASB12 to this compact \(\Omega\), obtaining one \(N\)
and constants \(B_x\ge0\) before \(F_0,\psi,u\) are given.

Now take arbitrary \(F_0,\psi\) satisfying the stated hypotheses.
The incoming SPR13 theorem, the exact original
\(\mathrm{SmoothRep.exists\_smooth\_repr}\) producer, constructs
one continuous automorphic representative \(\psi_0\) of this
same \(F_0\), smooth along all original products, such that
every \(R_l\psi_0\) is continuous and represents \(d_lF_0\).
HCU16's uniqueness of continuous representatives, or the
global-AE/full-support argument of RQC6, gives
\(\psi=\psi_0\) everywhere.
This is the only incoming smooth-representative endpoint here;
no subsequent multiplicity-map construction is used.

For each \(l\in\mathcal S\), ASB13 with \(z=xu\) and
\(M=B_x\mathcal N(u)^N\) gives
\[
       \int_\Omega|R_l\psi(xuo)|^2\,d\mu(o)
           \le B_x^2\mathcal N(u)^{2N}\|d_lF_0\|_2^2.
 \tag{ASB.41}
\]
The hypotheses needed for ASB6 now hold for \(\psi\), from
SPR13 on the same fixed level.  Substituting (ASB.41) into
(ASB.19), and using
\(\sum_l a_l^2\le(\sum_l a_l)^2\) for \(a_l\ge0\), gives
\[
 |\psi(xu)|^2
 \le\left(\sqrt{C_0}\,B_x\,\mathcal N(u)^N
                           \sum_{l\in\mathcal S}\|d_lF_0\|_2\right)^2.
 \tag{ASB.42}
\]
The quantity in parentheses is nonnegative, so taking square
roots proves (ASB.40) with \(C_x=\sqrt{C_0}B_x\).
The same proof works for rank zero, using the exponent-zero
convention already checked in ASB12.

The order of choices is explicit:
\[
       (D,\theta,V,\mu,\mathcal F,\rho,U)
       \ \longrightarrow\ (N,\mathcal S)\
       \longrightarrow\ x,C_x\
       \longrightarrow\ F_0,\psi,u.
 \tag{ASB.43}
\]
In particular neither \(N\), the list, nor \(C_x\) is chosen
again for a new vector or representative.  Applying the same
statement to a smooth derivative vector fixed by the same
level leaves \(N,\mathcal S\) unchanged.
It does not assert an exponent uniform over a different
previously chosen level, target or parameter-dependent datum.
% END ASB14

\paragraph{Application boundary.}
ASB3--ASB4 supply the actual chart and scalar integration
comparison used before the smooth-representative producer.
ASB14 consumes only its output SPR13.  The later construction
of an original \(\mathrm{MMap}\) may consume ASB14, so the
dependency order is acyclic.
All claims preserve the original group, \(\mathrm{oneP}\),
ordered derivatives, \(\mathrm{normInf}\), Haar measure and
quotient measure.  Initial boundedness, the actual target
\(\mathrm{FinPair}\), the selected Bundle and its readings
and rational tests remain separate original producer obligations.
